\section{Introduction and purpose}\label{sec:intro}
\status{CLOSED narrative material; owner voice review pending}
The objective of the FT1536 development is a security argument whose
objects can be traced to a specified implementation. Such an argument
requires more than a reduction between ideal games: the key law,
sampler law, randomness interface and byte-level acceptance predicate
must refer to the same experiment. The present artifact contains a
conditional classical reduction and several source-bound components;
their final implementation-level composition is unfinished.\src{C01}

This paper develops three parts of that argument. First, it presents the
conditional reduction and the computed constants of a pointwise sampler
comparison. Second, it records negative results that rule out shortcuts
in that comparison. Third, it describes the source-binding and review
process, including defects in both C code and formal models that the
process exposed. The artifact is part of the FT1536 evidence
archive~\cite{ft1536artifact}.\src{C04,C06,C09,C11,C12}

% Mission paragraph signed by the owner 2026-10-07 ("idealna wersja do
% paperu i do serca") - EN text below, PL original recorded in
% notes/PAPER_WORK_STATE.md.
FT1536 exists because trust in signature software should rest on
arguments that anyone can inspect, not on the authority of whoever
shipped the binary. The project develops ternary-secret signature
profiles in the Falcon lineage --- freely chosen, freely auditable,
freely usable --- to make independently inspectable cryptographic tools
available to people who wish to choose and evaluate such tools
themselves. That purpose requires publishing limitations and failed
arguments alongside successful proofs, binding claims to the exact code
that runs, and preserving the attribution of the Falcon Project and
Thomas Pornin. Accessibility is a design aim; it is not evidence for a
security claim.

Section~\ref{sec:main} distinguishes the target assembly from the
theorems already present. Section~\ref{sec:sampler} states the negative
results as theorems. Sections~\ref{sec:security} and~\ref{sec:limits}
give the security-accounting convention and the unresolved obligations
in the body of the paper.
The end-to-end assembly and its in-flight obligations are described in
Section~\ref{sec:main}.
